\subsection{Caché LRU de bloques descomprimidos}

El segundo componente clave del desarrollo fue la incorporación de una caché de bloques descomprimidos. Su objetivo es mantener en memoria de trabajo
los bloques usados recientemente, de modo que múltiples operaciones sucesivas puedan reutilizarlos sin repetir el proceso de descompresión.

La política de reemplazo adoptada es LRU. Cada acceso actualiza la marca de uso reciente del bloque, y cuando la caché se llena se selecciona para
expulsión el bloque con el acceso más antiguo. Si el bloque expulsado fue modificado, se marca como sucio y se recomprime antes de salir de la caché;
si no hubo modificaciones, puede descartarse sin costo adicional de escritura.

Desde el punto de vista de la simulación, esta decisión es importante porque muchas secuencias de compuertas reutilizan un subconjunto reducido del
estado. En tales casos, una caché pequeña pero bien administrada puede reducir de forma significativa la frecuencia de descompresiones y recompresiones. En la realización actual, la caché actúa además como mecanismo de \emph{write-back}: los bloques modificados permanecen como sucios hasta el vaciado, evitando escrituras comprimidas prematuras durante una secuencia corta de operaciones.
